\documentclass[11pt]{article}
\usepackage{fontspec}
\setmainfont{TeX Gyre Pagella}
\usepackage{amsmath,amssymb,amsthm}
\usepackage{geometry}
\geometry{margin=1.1in}
\usepackage{setspace}
\onehalfspacing
\usepackage{enumitem}
\usepackage[colorlinks=true,linkcolor=black,citecolor=black,urlcolor=black]{hyperref}

\newtheorem{definition}{Definition}
\newtheorem{proposition}{Proposition}
\newtheorem{corollary}{Corollary}
\theoremstyle{remark}
\newtheorem{remark}{Remark}
\newtheorem{example}{Example}

\newcommand{\R}{\mathbb{R}}
\newcommand{\N}{\mathcal{N}}
\newcommand{\id}{\operatorname{id}}
\newcommand{\rank}{\operatorname{rank}}
\newcommand{\E}{\mathbb{E}}
\newcommand{\tv}{\mathrm{TV}}

\title{Holocoordinate Systems:\\ Multidifference Observation, Transition Structure,\\ and Residual Growth in Enriched Atlases}
\author{Flyxion\\ \small Independent Researcher}
\date{September 2026}

\begin{document}
\maketitle

\begin{abstract}
\noindent A coordinate is ordinarily a number, or a vector of numbers, attached to a state. This essay studies a richer object, the \emph{holocoordinate system}: an atlas of perspectives in which each chart carries, besides its local coordinate, a relational signature linking the state to the wider field, translation maps into other perspectives, and a record of the discrepancies the chart fails to explain. Four properties of such a system are formalized and shown to be logically independent: \emph{multidifference completeness}, the requirement that the blind spaces of the charts intersect trivially, characterized through the kernel of the joint differential and quantified by a frame operator; \emph{holography}, the injectivity of a signature on the distinctions to be preserved; \emph{transition consistency}, the cocycle condition on translations; and \emph{holonomy}, the accumulated discrepancy of translation around closed loops of perspectives. Quantitative bounds are given for error amplification under ill-conditioned observation, for cocycle defects and holonomy in terms of pairwise discrepancies, and for the decrease of residuals when a chart is enlarged by a new coordinate, with a martingale limit for the growth rule. An application section treats a tracer-replication instrument as a chart, deriving its blind space, the gauge character of its holonomy under parameter redundancy, and criteria for transporting a fitted coordinate to another perspective. The results are conditional on smooth structure that is not assumed to be available for experiential data, and the essay states where that condition bites.
\end{abstract}

\section{Introduction}

A coordinate chart on a manifold assigns to each point of an open set a tuple of real numbers, and is the basic device by which a space is made available to calculation. In applied settings the device is often used with less care than in geometry. A chart is taken to have captured a state when its numbers have been recorded, without asking which changes of state leave the numbers unchanged, whether the numbers can be translated into those produced by another chart, or what remains unexplained once the numbers are used. The consequences are familiar from instrument-based measurement: parameters fitted within a model are read as properties of the phenomenon, when they are first of all coordinates of the model.

This essay develops a vocabulary in which those questions have exact answers. The starting point is a state space $X$, assumed to carry the structure of a smooth manifold where derivatives are used, and a family of \emph{perspectives}, each an observation map from a domain in $X$ into a feature space. Three demands are placed on the family.

\begin{enumerate}[label=(\roman*)]
\item \emph{Discrimination.} Every nonzero local difference in $X$ should register in at least one perspective, even if no single perspective registers all of them.
\item \emph{Translation.} The perspectives should retain the means of moving representations between one another, so that agreement, disagreement, and path dependence among them are measurable.
\item \emph{Accountability.} Each perspective should record what it fails to reproduce, so that persistent failure can motivate an enlarged chart rather than being absorbed as noise.
\end{enumerate}

Sections 3 through 6 treat these demands in turn, with a fourth section on signatures that connect a state to the wider field. Section 7 assembles the definitions into an enriched atlas and proves that its component properties are independent, and Section 8 applies the vocabulary to a tracer-replication instrument. The interest of the construction is partly that independence: a system can discriminate without reconstructing, reconstruct without translating consistently, and translate consistently on every overlap while carrying global holonomy. Keeping the properties separate prevents any one of them from serving as an all-purpose warrant for a coordinate system, in particular the word ``holographic''.

The motivating problem, that of coordinates supplied by an instrument whose parameterization may not span the phenomenon, has been discussed in an informal setting elsewhere \cite{flyxion}. What is offered here is the mathematics.

\section{Setting}

Let $X$ be a set of states. Where derivatives are used, $X$ is a smooth manifold of dimension $m$; where a probability structure is used, $X$ carries a probability measure $\mu$. A \emph{perspective} is a triple $\alpha=(U_\alpha,V_\alpha,\phi_\alpha)$ consisting of a domain $U_\alpha\subseteq X$, a feature space $V_\alpha$, and an observation map
\[
\phi_\alpha:U_\alpha\longrightarrow V_\alpha .
\]
In Sections 3 and 4, $V_\alpha$ is a finite-dimensional real inner product space, $U_\alpha$ is open, and $\phi_\alpha$ is $C^1$. In Sections 5 and 7, $V_\alpha$ need only be a metric space $(V_\alpha,d_\alpha)$. Different perspectives need not share dimension, variables, or units.

On an overlap $U_{\alpha\beta}=U_\alpha\cap U_\beta$, a \emph{transition map} is a map $T_{\beta\alpha}:V_\alpha\to V_\beta$ intended to satisfy $\phi_\beta=T_{\beta\alpha}\circ\phi_\alpha$ on $U_{\alpha\beta}$. An atlas in the classical sense also requires
\[
T_{\alpha\alpha}=\id,\qquad T_{\alpha\beta}=T_{\beta\alpha}^{-1},\qquad T_{\gamma\beta}\circ T_{\beta\alpha}=T_{\gamma\alpha},
\]
the last being the \emph{cocycle condition}: translating from $\alpha$ to $\gamma$ directly agrees with translating through $\beta$. For observational data the requirements will typically hold only approximately, and the departures are the objects of study.

\section{Blind spaces and multidifference completeness}

Fix $x\in X$ and let $A_x=\{\alpha:x\in U_\alpha\}$, assumed finite.

\begin{definition}[Blind space]
The \emph{blind space} of perspective $\alpha$ at $x$ is
\[
\N_\alpha(x)=\ker D\phi_\alpha(x)\subseteq T_xX .
\]
The \emph{joint blind space} at $x$ is
\[
\N(x)=\bigcap_{\alpha\in A_x}\N_\alpha(x)=\ker D\Phi_x,
\qquad
D\Phi_x(v)=\bigl(D\phi_\alpha(x)v\bigr)_{\alpha\in A_x}.
\]
The family is \emph{multidifference-complete at $x$} when $\N(x)=\{0\}$. In the vocabulary of differential topology this says that the joint observation map $\Phi$ is an immersion at $x$, or is jointly infinitesimally injective.
\end{definition}

A tangent direction $v$ lies in $\N_\alpha(x)$ when motion along $v$ produces no first-order change in perspective $\alpha$. This is the precise sense in which a chart has ``no coordinate'' for a direction: not that no coordinate exists anywhere, but that this observation map is locally insensitive to it. Completeness asks that the insensitivities of different perspectives do not overlap.

\begin{proposition}[Completeness is open and implies local injectivity]
The set of points at which the family is multidifference-complete is open. If the family is complete at $x$, there is a neighborhood $W$ of $x$ on which the joint observation map $\Phi=(\phi_\alpha)_{\alpha\in A_x}$ is injective, and $\Phi|_W$ is an embedding into $\bigoplus_{\alpha\in A_x}V_\alpha$.
\end{proposition}

\begin{proof}
Since the $U_\alpha$ are open and $A_x$ is finite, $A_y=A_x$ for $y$ near $x$. Injectivity of the linear map $D\Phi_y$ is equivalent to the nonvanishing of some $m\times m$ minor of its matrix in local coordinates, and this is an open condition because the entries depend continuously on $y$. Completeness at $x$ says that $\Phi$ is an immersion at $x$, and an immersion is locally an embedding by the rank theorem \cite{lee}.
\end{proof}

\begin{remark}
The converse fails. On $X=\R$ the map $\phi(x)=x^3$ is globally injective while $D\phi(0)=0$. First-order blindness at a point indicates degeneracy of the observation at that order, a direction that is ``technically observable but nearly invisible'', not necessarily a failure to distinguish states.
\end{remark}

To quantify how well the perspectives discriminate, equip $T_xX$ with an inner product $g_x$ and each $V_\alpha$ with its inner product $h_\alpha$. Define the \emph{frame operator}
\[
S_x=\sum_{\alpha\in A_x}D\phi_\alpha(x)^{*}D\phi_\alpha(x):T_xX\to T_xX,
\]
where the adjoint is taken with respect to $g_x$ and $h_\alpha$. In the vocabulary of frame theory \cite{christensen}, $S_x$ is the frame operator of the family of linear maps $\{D\phi_\alpha(x)\}$.

\begin{proposition}[Frame characterization]
The operator $S_x$ is self-adjoint and positive semidefinite, with $\langle S_xv,v\rangle=\sum_{\alpha}\|D\phi_\alpha(x)v\|^2$ and $\ker S_x=\N(x)$. Hence the family is complete at $x$ if and only if $S_x$ is positive definite. If $A(x)=\lambda_{\min}(S_x)$ and $B(x)=\lambda_{\max}(S_x)$, then
\[
A(x)\|v\|^2\;\le\;\sum_{\alpha\in A_x}\|D\phi_\alpha(x)v\|^2\;\le\;B(x)\|v\|^2
\qquad(v\in T_xX),
\]
and the stacked Jacobian $J(x)$ has extreme singular values $\sigma_{\min}=\sqrt{A}$, $\sigma_{\max}=\sqrt{B}$. Consequently completeness is equivalent to $\rank J(x)=m$.
\end{proposition}

\begin{proof}
The identity $\langle S_xv,v\rangle=\sum_\alpha\|D\phi_\alpha(x)v\|^2$ follows from the definition of the adjoint. If $S_xv=0$ then this sum vanishes, so $v\in\N(x)$; conversely $v\in\N(x)$ makes each term zero, so $\langle S_xv,v\rangle=0$, and positive semidefiniteness gives $S_xv=0$. The bounds are the Rayleigh bounds for $S_x$. Finally $S_x=J(x)^{*}J(x)$, so the eigenvalues of $S_x$ are the squares of the singular values of $J(x)$.
\end{proof}

Define the \emph{observational condition number}
\[
\kappa_{\mathrm{obs}}(x)=\frac{\sigma_{\max}(J(x))}{\sigma_{\min}(J(x))}=\sqrt{\frac{B(x)}{A(x)}},
\]
with the convention $\kappa_{\mathrm{obs}}=\infty$ when the family is incomplete.

\begin{remark}
The value of $\kappa_{\mathrm{obs}}$ depends on the chosen inner products $g_x$ and $h_\alpha$: rescaling the feature space of one perspective changes the weight with which its differential enters $S_x$. Completeness is invariant under such choices, since it depends only on kernels, but conditioning is not. It should be reported together with the weights, which encode how much each perspective's variables are trusted relative to the others.
\end{remark}

\begin{proposition}[Error amplification]
Suppose the family is complete at $x$ and one observes $y=J(x)v+\eta$ with noise $\eta$, estimating the difference by least squares, $\hat v=J(x)^{+}y$. Then
\[
\|\hat v-v\|\le\frac{\|\eta\|}{\sigma_{\min}(J(x))},
\qquad
\frac{\|\hat v-v\|}{\|v\|}\le\kappa_{\mathrm{obs}}(x)\,\frac{\|\eta\|}{\|J(x)v\|}\quad(v\neq0),
\]
and the first bound is attained for some direction of $\eta$.
\end{proposition}

\begin{proof}
Since $J$ has full column rank, $J^{+}=(J^{*}J)^{-1}J^{*}$ and $J^{+}J=\id$, so $\hat v-v=J^{+}\eta$. The operator norm of $J^{+}$ is $1/\sigma_{\min}(J)$ \cite{golub}, giving the first bound with equality when $\eta$ lies along the left singular vector for $\sigma_{\min}$. For the second, $\|J v\|\le\sigma_{\max}\|v\|$ gives $\|v\|\ge\|Jv\|/\sigma_{\max}$; dividing the first bound by $\|v\|$ yields the claim.
\end{proof}

An ill-conditioned family is therefore one in which some difference is observable in principle and amplified into large uncertainty in practice.

\begin{corollary}[Dimension count]\label{cor:dim}
(i) $\dim\N_\alpha(x)\ge m-\dim V_\alpha$; a perspective with $\dim V_\alpha<m$ therefore has a nonzero blind space. (ii) If the family is complete at $x$, then $\sum_{\alpha\in A_x}\dim V_\alpha\ge m$. (iii) If $\phi_\alpha$ is locally constant near $x$, then $\N_\alpha(x)=T_xX$.
\end{corollary}

\begin{proof}
(i) By rank-nullity, $\dim\ker D\phi_\alpha(x)=m-\rank D\phi_\alpha(x)\ge m-\dim V_\alpha$. (ii) Injectivity of $D\Phi_x:T_xX\to\bigoplus_\alpha V_\alpha$ requires $m\le\sum_\alpha\dim V_\alpha$. (iii) A locally constant map has zero differential.
\end{proof}

Part (iii) applies in particular to perspectives with discrete output, such as a set of category labels, which are locally constant away from category boundaries. First-order completeness cannot be obtained from them.

\section{Relational signatures and reconstruction}

The ``holo'' component of a holocoordinate is a signature that locates a state through its relations to the wider field rather than through independent axes. Let $(X,d)$ be a metric space. A \emph{landmark signature} with landmarks $\ell_1,\dots,\ell_N\in X$ is
\[
\sigma(x)=\bigl(d(x,\ell_1),\dots,d(x,\ell_N)\bigr).
\]
More generally a perspective $\alpha$ may carry a kernel signature $\sigma_\alpha(x)=\int_XK_\alpha(x,y)\Psi(y)\,d\mu(y)$. Such a signature is a compressed description of how the wider field appears from $x$. It is not, without further argument, holographic in any strong sense.

\begin{definition}[Holography]
Let $\sim$ be an equivalence relation on $X$ recording the distinctions one wishes to preserve. A signature $\sigma_\alpha$ is \emph{strongly holographic} (with respect to $\sim$) if
\[
\sigma_\alpha(x)=\sigma_\alpha(y)\ \Longrightarrow\ x\sim y .
\]
A family $\{\sigma_\alpha\}_{\alpha\in A}$ is \emph{jointly holographic} if the joint signature $\Sigma(x)=(\sigma_\alpha(x))_{\alpha\in A}$ has the same property.
\end{definition}

Strong holography is exactly the existence of a reconstruction: a map $G_\alpha$ with $G_\alpha(\sigma_\alpha(x))=[x]$ exists if and only if the implication above holds, and likewise jointly. The following elementary facts separate the two notions.

\begin{proposition}
(i) If some $\sigma_\alpha$ is strongly holographic, the family is jointly holographic. (ii) The converse fails: on $X=\R^2$ with $\sigma_1(x)=x_1$ and $\sigma_2(x)=x_2$ and $\sim$ equality, the family is jointly holographic while neither signature is.
\end{proposition}

\begin{proof}
(i) If $\Sigma(x)=\Sigma(y)$ then in particular $\sigma_\alpha(x)=\sigma_\alpha(y)$, so $x\sim y$. (ii) $\Sigma(x)=x$ is injective, while $\sigma_1$ identifies $(x_1,0)$ with $(x_1,1)$.
\end{proof}

In Euclidean space a landmark signature is holographic once enough landmarks are used, and the count is sharp.

\begin{proposition}[Trilateration]
Let $X=\R^m$ with the Euclidean distance and let $\ell_0,\dots,\ell_m$ be affinely independent. Then the landmark signature $\sigma(x)=(\|x-\ell_i\|)_{i=0}^m$ is injective. Fewer than $m+1$ landmarks do not suffice in general.
\end{proposition}

\begin{proof}
Expanding squared distances and subtracting the equation for $i=0$ from that for $i\ge1$ gives the linear system
\[
2(\ell_i-\ell_0)^{\top}x=\|\ell_i\|^2-\|\ell_0\|^2-\sigma_i^2+\sigma_0^2,\qquad i=1,\dots,m.
\]
The rows $(\ell_i-\ell_0)^{\top}$ are linearly independent by affine independence, so the matrix is invertible and $x$ is determined by $\sigma(x)$. With only $m$ landmarks the points $x$ and its reflection across the affine hull of the landmarks have the same signature.
\end{proof}

Injectivity of landmark maps is not automatic in other geometries. On a circle with the arc-length metric, a single landmark cannot separate a point from its reflection through the landmark's axis, while two landmarks that are neither equal nor antipodal do: two distinct reflections composing to a fixed point would force the axes to coincide. Whether a signature is holographic is a property to be verified for the space at hand.

\begin{remark}
If $X$ is compact Hausdorff and $\Sigma$ is continuous and injective, then $\Sigma$ is a homeomorphism onto its image and the reconstruction map is continuous. Without compactness a set-theoretic reconstruction may exist and still be discontinuous.
\end{remark}

\section{Transition consistency and holonomy}

From here on $V_\alpha$ is a metric space and each $T_{\beta\alpha}:V_\alpha\to V_\beta$ is a globally defined map, assumed Lipschitz with constant $L_{\beta\alpha}$ where a bound is stated. The \emph{transition discrepancy} of a state $x\in U_{\alpha\beta}$ is
\[
\delta_{\alpha\beta}(x)=d_\beta\bigl(\phi_\beta(x),\,T_{\beta\alpha}(\phi_\alpha(x))\bigr).
\]
This is not merely error. Nonzero discrepancy may reflect distortion, information genuinely specific to one perspective, an inadequate transition map, or a missing coordinate; the framework records it without deciding among these. The \emph{cocycle defect} of a triple is
\[
\kappa_{\alpha\beta\gamma}(y)=d_\gamma\bigl(T_{\gamma\beta}T_{\beta\alpha}y,\;T_{\gamma\alpha}y\bigr),\qquad y\in V_\alpha .
\]

\begin{proposition}[Cocycle defect from discrepancies]
For $x\in U_\alpha\cap U_\beta\cap U_\gamma$ and $y=\phi_\alpha(x)$,
\[
\kappa_{\alpha\beta\gamma}(y)\;\le\;L_{\gamma\beta}\,\delta_{\alpha\beta}(x)+\delta_{\beta\gamma}(x)+\delta_{\alpha\gamma}(x).
\]
In particular, if all discrepancies vanish, the cocycle condition holds at $\phi_\alpha(x)$.
\end{proposition}

\begin{proof}
By the triangle inequality in $V_\gamma$,
\begin{align*}
\kappa_{\alpha\beta\gamma}(y)
&\le d_\gamma\bigl(T_{\gamma\beta}T_{\beta\alpha}y,\,T_{\gamma\beta}\phi_\beta(x)\bigr)
+d_\gamma\bigl(T_{\gamma\beta}\phi_\beta(x),\,\phi_\gamma(x)\bigr)
+d_\gamma\bigl(\phi_\gamma(x),\,T_{\gamma\alpha}y\bigr).
\end{align*}
The first term is at most $L_{\gamma\beta}\,d_\beta(T_{\beta\alpha}y,\phi_\beta(x))=L_{\gamma\beta}\delta_{\alpha\beta}(x)$, the second is $\delta_{\beta\gamma}(x)$, and the third is $\delta_{\alpha\gamma}(x)$.
\end{proof}

Now consider a loop of perspectives $\alpha_0\to\alpha_1\to\cdots\to\alpha_k=\alpha_0$ and write $T_i=T_{\alpha_{i+1}\alpha_i}$ with Lipschitz constants $L_i$. The \emph{holonomy map} of the loop is
\[
\mathcal H_\gamma=T_{k-1}\circ\cdots\circ T_1\circ T_0:V_{\alpha_0}\to V_{\alpha_0}.
\]
For a state $x$ lying in every $U_{\alpha_i}$, the \emph{holonomy discrepancy} is
\[
\Omega_\gamma(x)=d_{\alpha_0}\bigl(\phi_{\alpha_0}(x),\,\mathcal H_\gamma(\phi_{\alpha_0}(x))\bigr).
\]

\begin{proposition}[Holonomy bound]
For $x\in\bigcap_{i}U_{\alpha_i}$, writing $\delta_j=\delta_{\alpha_j\alpha_{j+1}}(x)$,
\[
\Omega_\gamma(x)\;\le\;\sum_{j=0}^{k-1}\Bigl(\prod_{l=j+1}^{k-1}L_l\Bigr)\delta_j .
\]
In particular $\Omega_\gamma(x)>0$ implies $\delta_j(x)>0$ for some $j$, and a loop over a common state has trivial holonomy discrepancy whenever the discrepancies vanish.
\end{proposition}

\begin{proof}
Let $z_i=\phi_{\alpha_i}(x)$ and define $y_0=z_0$, $y_{i+1}=T_i(y_i)$, so that $y_k=\mathcal H_\gamma(z_0)$ and $z_k=z_0$. We show by induction that
\[
d_{\alpha_i}(y_i,z_i)\le\sum_{j<i}\Bigl(\prod_{l=j+1}^{i-1}L_l\Bigr)\delta_j .
\]
For $i=0$ both sides are zero. Given the bound at $i$,
\[
d(y_{i+1},z_{i+1})\le d(T_iy_i,T_iz_i)+d(T_iz_i,z_{i+1})\le L_i\,d(y_i,z_i)+\delta_i,
\]
and substituting the inductive bound gives the bound at $i+1$. Taking $i=k$ and using $\Omega_\gamma(x)=d_{\alpha_0}(z_k,y_k)$ completes the proof.
\end{proof}

The bound runs one way. Discrepancies may cancel around a loop, so vanishing holonomy does not imply vanishing discrepancy: holonomy detects the part of the discrepancy that survives transport, and is a coarser instrument than the pairwise discrepancies themselves. Around a triangle the relation to the cocycle condition is exact.

\begin{proposition}[Triangle holonomy]
Suppose $T_{\alpha\gamma}=T_{\gamma\alpha}^{-1}$ with $T_{\gamma\alpha}$ bijective, and let
\[
H=T_{\alpha\gamma}T_{\gamma\beta}T_{\beta\alpha}.
\]
Then for $y\in V_\alpha$, $H(y)=y$ if and only if $\kappa_{\alpha\beta\gamma}(y)=0$.
\end{proposition}

\begin{proof}
$H(y)=y$ holds if and only if $T_{\gamma\alpha}^{-1}(T_{\gamma\beta}T_{\beta\alpha}y)=y$, that is, if and only if $T_{\gamma\beta}T_{\beta\alpha}y=T_{\gamma\alpha}y$.
\end{proof}

Loops that are not products of triangles are not controlled by the cocycle conditions. The following example shows that exact charts can carry nontrivial holonomy even though every cocycle condition holds.

\begin{example}[Winding on the circle]\label{ex:winding}
Let $X=\R/2\pi\mathbb Z$ and choose $0<\varepsilon<\pi/3$. Cover $X$ by three arcs $U_1=(-\varepsilon,\tfrac{2\pi}{3}+\varepsilon)$, $U_2=(\tfrac{2\pi}{3}-\varepsilon,\tfrac{4\pi}{3}+\varepsilon)$, $U_3=(\tfrac{4\pi}{3}-\varepsilon,2\pi+\varepsilon)$, with $\phi_\alpha:U_\alpha\to\R$ the angle read in the stated interval. The overlaps $U_{12}$, $U_{23}$, $U_{13}$ are nonempty and there are no triple overlaps. Take $T_{21}=T_{32}=\id_\R$ and $T_{31}(y)=y+2\pi$, which agree with $\phi_\beta\circ\phi_\alpha^{-1}$ on the overlaps, the last because a point near $0$ has angle in $(-\varepsilon,\varepsilon)$ in chart 1 and in $(2\pi-\varepsilon,2\pi+\varepsilon)$ in chart 3. Every chart is an exact coordinate and every cocycle condition is vacuous, yet for the loop $1\to2\to3\to1$,
\[
\mathcal H=T_{13}\circ T_{32}\circ T_{21}=T_{31}^{-1},\qquad \mathcal H(y)=y-2\pi\neq y .
\]
The holonomy is the winding number of the circle, invisible to every chart and to every triple of charts.
\end{example}

\begin{remark}[Cohomological classification]
Suppose the transition maps are invertible, take values in a group $G$ acting on a common feature space (for instance translations or affine maps), and satisfy the cocycle condition on every nonempty triple overlap. Then $\{T_{\beta\alpha}\}$ is a \v{C}ech $1$-cocycle on the nerve $\mathcal{K}$ of the cover with values in $G$. Its class in $\check H^1(\mathcal K;G)$ vanishes if and only if there are $g_\alpha\in G$ with $T_{\beta\alpha}=g_\beta g_\alpha^{-1}$, that is, if and only if the charts can be renormalized by $\phi_\alpha\mapsto g_\alpha^{-1}\phi_\alpha$ so that every transition map becomes the identity. For connected $\mathcal K$ the classes correspond to $\mathrm{Hom}(\pi_1(\mathcal K),G)/G$, the class being the holonomy representation, so holonomy around every loop is trivial exactly when the class is trivial. When $\mathcal K$ is simply connected the class vanishes for every $G$, which is why the triangle conditions then control everything. In Example~\ref{ex:winding} the nerve is the boundary of a triangle with $\pi_1=\mathbb Z$, $G=(\R,+)$, and the class is $2\pi\in\check H^1(\mathcal K;\R)=\R$. For nonabelian $G$ the set $\check H^1$ is a pointed set rather than a group \cite{bott}. Transition maps that are merely Lipschitz, without a group structure, are not covered by this classification.
\end{remark}

\section{Residuals and the growth rule}

A perspective can also be held to account for what it fails to reproduce. Let $P_q:X\to\R$ be a family of square-integrable \emph{probes} $q\in Q$ extracting properties from a state, and let $\phi_\alpha$ be a chart whose output is used to predict them. With squared error, the best decoder of probe $q$ from the chart is the conditional expectation, and the \emph{residual} is
\[
\varepsilon_{\alpha,q}(x)=P_q(x)-\E\bigl[P_q\mid\phi_\alpha\bigr](x),
\qquad
\varepsilon_\alpha=(\varepsilon_{\alpha,q})_{q\in Q}.
\]
Enlarging the chart by a measurable function $z$ of the state gives $\tilde\phi_\alpha=(\phi_\alpha,z)$ and the enlarged residual $\tilde\varepsilon_{\alpha,q}=P_q-\E[P_q\mid\phi_\alpha,z]$.

\begin{proposition}[Residual reduction]\label{prop:resid}
For each probe,
\[
\E\bigl[\tilde\varepsilon_{\alpha,q}^{\,2}\bigr]
=\E\bigl[\varepsilon_{\alpha,q}^{\,2}\bigr]
-\E\Bigl[\bigl(\E[P_q\mid\phi_\alpha,z]-\E[P_q\mid\phi_\alpha]\bigr)^2\Bigr].
\]
Hence enlargement never increases the mean squared residual, and decreases it strictly if and only if $\E[P_q\mid\phi_\alpha,z]\neq\E[P_q\mid\phi_\alpha]$ on a set of positive measure, that is, if and only if the residual is \emph{structured} with respect to $z$.
\end{proposition}

\begin{proof}
Write $P=P_q$, $F=\sigma(\phi_\alpha)$, $\tilde F=\sigma(\phi_\alpha,z)\supseteq F$. Then
\[
P-\E[P\mid F]=\bigl(P-\E[P\mid\tilde F]\bigr)+\bigl(\E[P\mid\tilde F]-\E[P\mid F]\bigr).
\]
The two summands are orthogonal in $L^2$: the first is orthogonal to every $\tilde F$-measurable square-integrable function, and the second is $\tilde F$-measurable. Taking second moments gives the identity.
\end{proof}

This yields the growth rule
\[
\begin{aligned}
&\text{observe}\to\text{translate}\to\text{detect residual}\\
&\qquad\to\text{differentiate residual}\to\text{add coordinate}
\end{aligned}
\]
a quantitative form. A residual that is unstructured, in the sense that $\E[\varepsilon_{\alpha,q}\mid\phi_\alpha,z]=0$ for every candidate $z$, cannot be reduced by any enlargement.

\begin{corollary}[Limit of the growth rule]\label{cor:growth}
Let $z_1,z_2,\dots$ be successively added coordinates and $\mathcal F_n=\sigma(\phi_\alpha,z_1,\dots,z_n)$, $\mathcal F_\infty=\sigma(\bigcup_n\mathcal F_n)$. Then the mean squared residual is nonincreasing in $n$ and converges to $\E[(P_q-\E[P_q\mid\mathcal F_\infty])^2]$. The limit is zero if and only if $P_q$ is $\mathcal F_\infty$-measurable almost surely.
\end{corollary}

\begin{proof}
Monotonicity is the previous proposition. Convergence of $\E[P_q\mid\mathcal F_n]$ to $\E[P_q\mid\mathcal F_\infty]$ in $L^2$ is L\'evy's upward theorem \cite{williams}, and the last statement follows because a function equals its conditional expectation on $\mathcal F_\infty$ exactly when it is $\mathcal F_\infty$-measurable.
\end{proof}

\begin{remark}
The statement is population-level. With finite data, enlargement can fit noise, and detecting a structured residual requires held-out probes; no claim is made here about estimation.
\end{remark}

\section{Enriched atlases and the independence of their properties}

\begin{definition}[Holocoordinate system]
A \emph{holocoordinate system} on $X$ is an enriched atlas
\[
\mathfrak H=\bigl\{(U_\alpha,\phi_\alpha,\sigma_\alpha,T_{\beta\alpha},\varepsilon_\alpha)\bigr\}_{\alpha\in A},
\]
in which $\phi_\alpha$ is a local coordinate, $\sigma_\alpha$ is a relational signature, $T_{\beta\alpha}$ transports representations between perspectives, and $\varepsilon_\alpha$ records the residuals of the chart. For a designated set $A_0\subseteq A$ of perspectives, $\mathfrak H$ is
\begin{itemize}[nosep]
\item[\textbf{(M)}] \emph{multidifference-complete at $x$} if $\N(x)=\{0\}$;
\item[\textbf{(H)}] \emph{strongly holographic} if $\sigma_\alpha$ is strongly holographic for every $\alpha\in A_0$;
\item[\textbf{(C)}] \emph{transition-consistent} if $\delta_{\alpha\beta}\equiv0$ on every overlap and all cocycle conditions hold;
\item[\textbf{(V)}] \emph{holonomy-trivial} if every holonomy map is the identity on its domain.
\end{itemize}
\end{definition}

\begin{proposition}[Independence]
None of the following implications holds in general.
\begin{enumerate}[label=(\alph*),nosep]
\item \textbf{(M)} does not imply \textbf{(H)}.
\item \textbf{(H)} does not imply \textbf{(C)}.
\item \textbf{(C)} does not imply \textbf{(V)}.
\item \textbf{(V)} does not imply \textbf{(C)}.
\item \textbf{(C)} does not imply \textbf{(M)}, and \textbf{(H)} does not imply \textbf{(M)}.
\end{enumerate}
\end{proposition}

\begin{proof}
(a) The family $\phi_1=x_1$, $\phi_2=x_2$ on $\R^2$ with $\sigma_\alpha=\phi_\alpha$ is complete, since $D\Phi=\id$, while neither signature is strongly holographic. (b) On $X=\R$ take three charts with $\phi_\alpha=\id$, $\sigma_\alpha=\id$ (injective), $T_{21}=T_{32}=\id$ and $T_{31}(y)=y+1$. Then $\kappa_{123}\equiv1$ and $\delta_{13}\equiv1$. (c) Example~\ref{ex:winding}: every chart is exact and all cocycle conditions hold vacuously, but the holonomy is a nonzero translation. (d) Take two charts $\phi_1=\phi_2=\id$ on $\R$ with $T_{21}(y)=y+1$ and $T_{12}=T_{21}^{-1}$. The loop $1\to2\to1$ is the identity, yet $\delta_{12}\equiv1$. (e) A single chart $\phi(x)=x_1$ on $\R^2$ satisfies \textbf{(C)} and \textbf{(V)} vacuously and fails \textbf{(M)}; a constant chart $\phi\equiv0$ with $\sigma=\id$ satisfies \textbf{(H)} and fails \textbf{(M)}.
\end{proof}

The independence is the point of the construction. A system can observe every difference without reconstructing the state (a), reconstruct without translating consistently (b), be exact on every overlap while carrying global holonomy (c), close every loop without being consistent (d), and be consistent or reconstructive while blind to some directions (e). Each property is separately testable, and a claim that a system is ``holographic'' is accordingly a claim about signature injectivity and nothing else.

\section{Application: the tracer replication chart}

The preceding vocabulary applies to any instrument that converts a compressed parameter vector into an exposed rendering which a participant compares with an experience. The Tracer Replication Tool discussed in \cite{flyxion} and \cite{qri_note} is such an instrument. This section uses it as an illustration and makes no claim about any data collected with it.

\paragraph{Modeling assumptions.} (A1) The experiences of interest form a smooth $m$-manifold $X$. (A2) Comparison is mediated by a finite-dimensional Euclidean feature space $W$, with a smooth map $\iota:X\to W$ giving the features of an experience and a smooth renderer $R:\Theta\to W$, $\Theta\subseteq\R^p$ open, giving the features of the pattern produced by a parameter vector (persistence, intensity, decay, strobe frequency, replay frequency, pulse, envelope parameters, color inversion). (A3) The comparison functional is squared distance, $D_e(w)=\|\iota(e)-w\|^2$. These are strong simplifications. The subjective comparison of the participant need not be a metric, and (A1) is the smoothness hypothesis flagged in Section~\ref{sec:scope}.

Suppose first that $R$ is an embedding of $\Theta$ onto a submanifold $M=R(\Theta)\subseteq W$ of dimension $p$, and that $\mathcal T$ is a tubular neighborhood of $M$ with smooth nearest-point retraction $\pi:\mathcal T\to M$. For $\iota(e)\in\mathcal T$ the fitted parameter is unique, and the \emph{tool chart} is
\[
\phi_\tau(e)=R^{-1}\bigl(\pi(\iota(e))\bigr)\in\Theta .
\]

\begin{proposition}[Blind space of the tool chart]
If $\iota$ is an immersion at $e$, then
\[
\N_\tau(e)=\bigl\{v\in T_eX:\ D\iota(e)v\in N_{\pi(\iota(e))}M\bigr\},
\]
where $N_qM$ is the normal space of $M$ at $q$. Consequently $\dim\N_\tau(e)\ge m-p$.
\end{proposition}

\begin{proof}
By the tubular neighborhood theorem $\pi$ is a submersion whose fibers are the normal disks, so $\ker D\pi_w=N_{\pi(w)}M$. Since $R^{-1}:M\to\Theta$ is a diffeomorphism, $D\phi_\tau(e)=D(R^{-1})\circ D\pi\circ D\iota(e)$ has the kernel stated. Because $D\iota(e)$ is injective, $\dim\N_\tau(e)=\dim\bigl(D\iota(e)(T_eX)\cap N_qM\bigr)$, and two subspaces of $W$ of dimensions $m$ and $\dim W-p$ meet in dimension at least $m+(\dim W-p)-\dim W=m-p$.
\end{proof}

The blind space of the tool chart is thus the set of changes in experience whose perceptual image is normal to the model manifold: exactly the variation for which the parameterization has no coordinate. If $m>p$ the tool chart cannot be complete on its own (Corollary~\ref{cor:dim}), and completeness must come from other perspectives whose blind spaces meet the tool's only at zero.

\paragraph{From embedding to degeneracy.} Proposition 11 assumes that $R$ is an embedding. The behavior of the tool chart depends on how far $R$ departs from that, and three regimes should be distinguished.
\begin{enumerate}[label=(\arabic*),nosep]
\item \emph{Embedding.} The fit is unique near $\iota(e)$ and the blind space is given by Proposition 11.
\item \emph{Immersion that is not injective.} Locally $R$ is still an embedding, so on each small sheet Proposition 11 applies, but the sheets can overlap or wrap, and the fitted parameter is multivalued. The chart exists only on branches, and the blind space is a property of the branch.
\item \emph{Rank drop.} $DR(\theta)$ has a kernel, the fibers of $R$ have positive dimension, and the chart takes values in a quotient, as discussed below.
\end{enumerate}

\begin{proposition}[Branch dependence of the blind space]
Let $R$ be an immersion, let $\theta_0\neq\theta_1$ satisfy $R(\theta_0)=R(\theta_1)=\iota(e)=q$, and choose neighborhoods $\Theta_j\ni\theta_j$ on which $R$ is an embedding, with images $M_j$ admitting tubular retractions $\pi_j$ near $q$. Let $\phi^{(j)}(e')=(R|_{\Theta_j})^{-1}(\pi_j(\iota(e')))$ be the corresponding branch charts. If $\iota$ is an immersion at $e$, then
\[
\N^{(j)}(e)=D\iota(e)^{-1}\bigl(N_qM_j\bigr),\qquad j=0,1.
\]
In particular, if $T_qM_0=T_qM_1$ the two blind spaces coincide. In general they need not.
\end{proposition}

\begin{proof}
Apply Proposition 11 to each sheet. Equal tangent spaces have equal orthogonal complements in $W$, so the preimages coincide. For a failure of coincidence, take $X=W=\R^2$, $\iota=\id$, and two sheets crossing transversally at $q$: the blind spaces are the two distinct normal lines.
\end{proof}

\begin{example}[Monodromy in the tool setting]\label{ex:mono}
Let $\Theta=\R$, $W=\R^2$, $R(\theta)=(\cos\theta,\sin\theta)$, $X=W\setminus\{0\}$ and $\iota$ the inclusion, so $m=2$, $p=1$. Then $R$ is an immersion but not injective, $M$ is the unit circle, and the fitted parameter of an experience $e$ is its argument, determined only modulo $2\pi$. On each branch the blind space is the radial line, by Proposition 11. Continuing a branch chart along the loop $e(t)=(\cos t,\sin t)$, $t\in[0,2\pi]$, returns the parameter $\theta+2\pi$, so the holonomy is translation by $2\pi$, which is Example~\ref{ex:winding} realized inside the instrument. No continuous single-valued choice of parameter exists over the loop, since one would be a continuous lift of the argument map on the circle to the covering $\R\to S^1$. The obstruction disappears only if the chart is allowed to take values in the quotient $\Theta/{\sim_R}\cong S^1$, that is, only if the linear coordinate is given up.
\end{example}

\paragraph{Parameter redundancy.} Suppose instead that $R$ is not an immersion, so that distinct parameter vectors can produce indistinguishable outputs. Write $\theta\sim_R\theta'$ when $R(\theta)=R(\theta')$. The fitted parameter is then determined only up to $\sim_R$, and the tool chart takes values in the quotient $\Theta/{\sim_R}$. A function $f:\Theta\to\R$ is a well-defined function of the chart only if it respects the quotient.

\begin{proposition}[Well-defined coordinate functions]
A function $f:\Theta\to\R$, such as a fitted replay frequency, descends to a function on the chart values $\Theta/{\sim_R}$ if and only if $f$ is constant on the fibers of $R$. In particular, exhibiting $\theta,\theta'$ with $R(\theta)=R(\theta')$ and $f(\theta)\ne f(\theta')$ shows that $f(\theta^\ast)$ is not determined by what the tool distinguishes.
\end{proposition}

\begin{proof}
A function on a quotient is precisely a function constant on equivalence classes, and the classes of $\sim_R$ are the fibers of $R$.
\end{proof}

\begin{example}[Holonomy as gauge]
Let $s:R(\Theta)\to\Theta$ be a selection of one representative from each fiber, and consider the loop that renders a parameter vector, refits, and reads the parameters again: $\theta\mapsto R(\theta)\mapsto s(R(\theta))$. The holonomy discrepancy is $\Omega(\theta)=d(\theta,s(R(\theta)))$, which is positive exactly when $\theta$ is not the selected representative. It vanishes identically for the quotient chart. This holonomy is a coordinate artifact: it depends on the choice of $s$ and disappears in $\Theta/{\sim_R}$. It differs in kind from the winding of Example~\ref{ex:winding} and the monodromy of Example~\ref{ex:mono}, which cannot be removed by recoordinatizing the charts.
\end{example}

\paragraph{Categorical and graded reports.} Verbal categories such as \emph{trail}, \emph{strobe}, and \emph{replay} form a perspective with discrete output. By Corollary~\ref{cor:dim}(iii) such a perspective has full blind space wherever it is differentiable, and registers change only when a state crosses a category boundary. It can be handled by the metric machinery of Section~5 but contributes nothing to first-order completeness. Graded ratings on continuous scales do contribute, and a complete family in practice requires perspectives of that kind.

\paragraph{Transporting a fitted coordinate to a process.} A value of $\phi_\tau$ is a coordinate of the tool chart. To assert that it corresponds to a process in another substrate, say a feature vector $\phi_\nu$ of recorded neural activity, requires a transition map $T_{\nu\tau}$ and evidence about it. In the vocabulary of Sections 5 and 7 the assertion is supported to the extent that:
\begin{enumerate}[label=(\roman*),nosep]
\item the discrepancy $\delta_{\tau\nu}(e)$ is small on held-out states, that is, the map predicts the second perspective rather than being fitted to it;
\item cocycle defects $\kappa_{\tau\nu\gamma}$ are small with respect to further perspectives $\gamma$, such as psychophysical or behavioral charts, so that the correspondence is not particular to one pair;
\item the joint family is complete at the states in question, so that the agreement is not carried entirely by directions in which some chart is blind; and
\item the transported quantity is a well-defined function of the chart in the sense of the last proposition.
\end{enumerate}
A fitted replay frequency that satisfies none of these establishes that a replay at that frequency in the renderer resembles a remembered experience, and nothing further.

\paragraph{Contaminated growth.} Proposition~\ref{prop:resid} compares residuals under a fixed joint law of probes and chart outputs. If exposure to the tool changes how experiences are recalled and reported, the law after $k$ exposures, $\mu_k$, differs from $\mu_{k-1}$, and mean squared residuals computed under different laws are not comparable: a decrease may reflect changed reports rather than an added coordinate. The identity of Proposition~\ref{prop:resid} remains valid at each fixed $k$, and the growth rule remains meaningful if the probes are elicited, or the residuals evaluated, under a fixed reference law, for instance from reports collected before exposure. A candidate new coordinate $z$, such as a graded rating of spatial coherence of the trails that the tool does not parameterize, is supported when $\E[P_q\mid\phi_\tau,z]\ne\E[P_q\mid\phi_\tau]$ under that reference law, and not merely when participants report wanting a further control.

\section{Scope and limits}\label{sec:scope}

The results above are conditional on structure that must be argued for in any application.

First, the differential-geometric statements require that $X$ be a smooth manifold and the $\phi_\alpha$ differentiable. For state spaces of experience or of relational configuration this is a modeling hypothesis, not a finding. Where it fails, the algebraic content of Sections 5 through 7 survives, since it uses only metrics and Lipschitz maps, while Section 3 must be replaced by a discrete or set-theoretic analogue of distinguishability.

Second, the derivatives $D\phi_\alpha(x)$ are not observable from finitely many reports. In practice they would be approximated by finite differences over nearby states, and the completeness criterion becomes a statement about the numerical rank of an estimated Jacobian, with the stability caveats of Section 3.

Third, completeness at first order is a local and infinitesimal property. It does not guarantee global injectivity of the joint observation map, and it says nothing about whether the perspectives span the phenomenon rather than the instrument. A coordinate fitted within a family of charts is first of all a coordinate of that family; the growth rule of Section 6 is the mechanism by which persistent, structured failure can extend the family, and the limit in Corollary~\ref{cor:growth} states exactly which residual survives when the added coordinates are exhausted.

\begin{thebibliography}{9}
\bibitem{lee} Lee, J. M. (2013). \emph{Introduction to Smooth Manifolds} (2nd ed.). Springer.
\bibitem{christensen} Christensen, O. (2003). \emph{An Introduction to Frames and Riesz Bases}. Birkh\"auser.
\bibitem{golub} Golub, G. H., and Van Loan, C. F. (2013). \emph{Matrix Computations} (4th ed.). Johns Hopkins University Press.
\bibitem{bott} Bott, R., and Tu, L. W. (1982). \emph{Differential Forms in Algebraic Topology}. Springer.
\bibitem{williams} Williams, D. (1991). \emph{Probability with Martingales}. Cambridge University Press.
\bibitem{qri_note} Qualia Research Institute. \emph{Tracer Replication Tool} (Psychophysics Toolkit), and accompanying writeup by A.\ G\'omez Emilsson and colleagues. \url{https://qri.org}.
\bibitem{flyxion} Flyxion. \emph{Rehearsable Abstraction: Editable Expansion as a Regime of Implementation Visibility}. Unpublished essay, 2026.
\end{thebibliography}

\end{document}
